\section{Enforcement points}\label{sec:enforcement}
Part~2 defines how a decision is bound to a real operation: the adapter contract, the evidence that crosses each internal boundary, the tool executor, approvals and reservations, unknown outcomes, data-flow limits, and the MCP transport. Content-specific MCP operations are in Part~3.

\subsection{Adapter contract}
\begin{lstlisting}
// Logical design; generated DTOs are defined by the versioned API.
type EnforcementAdapter interface {
    Prepare(ctx context.Context, request AdapterRequest) (OperationPlan, error)
    Execute(ctx context.Context, plan OperationPlan, grant Grant) (Outcome, error)
}

type OperationPlan struct {
    ID              string
    RegistryVersion string
    RequiredChecks  []Check
    ArgumentsDigest string
    ResourceBindings []ResourceBinding
    IdempotencyKey  string
}
\end{lstlisting}
\req{ENF-001}{R1-B}{T-04}{Every adapter implements \code{Prepare} and \code{Execute}. \code{Prepare} resolves targets and validates arguments without performing the protected side effect. \code{Execute} verifies the grant, carries out the required obligations and dispatches only the prepared plan.}
\req{ENF-002}{R1-B}{T-04}{An adapter request is a validated typed protocol or domain variant. It \mustnot\ be an unrestricted \code{any} or an arbitrary attribute map, and each concrete adapter accepts only its declared request kind.}
\req{ENF-003}{R1-B}{T-01, T-04}{An adapter translates authentication evidence, selectors, operation parameters and response framing. It \mustnot\ fork business grants, duplicate policy combination, invent its own context format, turn an unknown policy result into a local allow, or perform a permissive check and then bypass the shared service.}
\req{ENF-004}{R1-B}{T-04}{Where target resolution is itself sensitive, it requires a separate discovery or stat action.}
\req{ENF-005}{R1-B}{T-04}{Every adapter declares its discovery, read, mutation, output-transformation, batching, streaming and obligation capabilities. An unknown capability fails closed. Public API clients receive redacted errors while internal audit retains diagnostic detail, and the resource-to-action mapping is versioned with policy and schema releases.}
\req{ENF-006}{R1-C}{T-04}{Database row or column filtering, filesystem mounts, message consumption and job scheduling are extension points that need real enforcement at their own access boundary. The first release does not compile expressions into SQL; a later optimizer must translate a declared subset, prove equivalence against per-resource decisions and never drop an unsupported predicate. A database adapter exposes a field only if it actually builds a restricted query or response.}

\subsection{Trust propagation}\label{sec:trust}
Authentication at an ingress is insufficient if a downstream service can be reached another way. Each internal call carries the evidence in Table~\ref{tab:hops}, and the receiver rejects a call that lacks it.
\begin{longtable}{@{}P{0.06\textwidth}P{0.20\textwidth}P{0.46\textwidth}P{0.10\textwidth}@{}}
\caption{Evidence that crosses each boundary}\label{tab:hops}\\
\toprule TB & Hop & Evidence the receiver verifies & Phase\\\midrule\endfirsthead
\toprule TB & Hop & Evidence the receiver verifies & Phase\\\midrule\endhead
TB1 & Client \arrowto\ adapter & Token or credential for this audience (\rr{ID-001}); session bound to actor, subject and delegation (\rr{MCP-007}) & R1\\
TB2 & Adapter \arrowto\ decision & Workload certificate bound to the adapter role (\rr{DEC-007}); references, not asserted facts (\rr{ID-003}) & R1-A\\
TB2 & Decision \arrowto\ adapter & Signed decision for this audience, operation and validity (\rr{DEC-008}) & R1-A\\
TB3 & MCP adapter \arrowto\ content service & Workload certificate of the \code{tools} role, plus signed decision or a delegation reference verified against authoritative state (\rr{TRU-003}) & R1-C\\
TB3 & Content service \arrowto\ broker & Workload certificate of the \code{storage} role plus a signed decision whose checks are exactly those of the matching operation mapping of Table~\ref{tab:keyops}, bound to the exact object, version, unit and recipient (\rr{TRU-005}) & R2\\
TB4 & Executor \arrowto\ connector & Target-specific restricted credential, never the model's credential (\rr{AGT-006}) & R1-B\\
TB4 & Executor \arrowto\ downstream & Idempotency key and state precondition where supported (\rr{AGT-006}, \rr{INV-005}) & R1-B\\
TB5 & Storage \arrowto\ blob backend & Verified workload binding on every call (\rr{BKD-001}) & R1-C\\
\bottomrule
\end{longtable}
\req{TRU-001}{R1-A}{T-03}{Each cross-process hop carries the evidence listed in Table~\ref{tab:hops}. A receiver that does not find it rejects the call.}
\req{TRU-002}{R1-A}{T-03, T-20}{Workload identity uses mTLS or an equivalent bound credential issued by the workload certificate authority chosen in decision D-04, bound to one role. Certificate verification is mandatory configuration, not an optional check.}
\reqn{TRU-003}{R1-A}{T-03}{The end user's authority crosses a hop only as a signed decision (\rr{DEC-008}) or as a delegation reference that the receiver verifies against authoritative state. Caller-asserted claims plus a workload identity are never sufficient.}
\reqn{TRU-004}{R1-A}{T-03}{A role accepts only signed decisions whose audience is that role, and the decision-signing key is registered separately from workload, issuer and encryption keys.}
\reqn{TRU-005}{R2}{T-11}{The broker releases key material only against a signed decision whose checks are exactly those of the matching operation mapping in Table~\ref{tab:keyops} and that is bound to the exact object, version, unit and recipient being released. The broker verifies the complete signed result, not a subset, whether or not the request also carries an mTLS identity. A compromised gateway therefore cannot obtain keys that no valid decision covers.}
\req{TRU-006}{R1-C}{T-03}{The MCP adapter calls the same content-service implementation as every other adapter and preserves the represented principal and delegation. It \mustnot\ replace the user with a broadly privileged connector account.}

\subsection{Tool executor}
The agent proposes an invocation. A trusted executor authenticates it, resolves the registered tool, validates arguments, computes the real operation plan, obtains authorization and performs the effect with credentials the model never sees.
\req{AGT-001}{R1-B}{T-04}{A registry record binds tool ID, version, schema digest, connector implementation, resource mapper, allowed destinations, required checks and supported obligations. Tool descriptions and annotations are untrusted unless supplied by the registry, and a label such as ``read-only'' does not establish behavior. A changed schema digest invalidates the registration until re-reviewed.}
\req{AGT-002}{R1-B}{T-04}{Hiding a tool from discovery improves usability but is not a security boundary. Every direct invocation is checked.}
\req{AGT-003}{R1-B}{T-04}{The executor performs these steps in order and \mustnot\ skip or reorder them. (1) Validate actor identity, represented subject, tenant, audience-bound session and delegation, and reject on-behalf-of headers supplied by the caller. (2) Resolve the tool and version and validate arguments against a strict schema, rejecting unknown fields where security depends on exact interpretation. (3) Map the request to concrete resources and actions (for example a submission, an account, a currency and amount, and a destination service), canonicalize them and compute the operation digest. (4) Evaluate every required check, resolve approval, reservation, destination restriction and other obligations before execution, and bind them to the exact plan. (5) At dispatch recheck expiry, generations, delegation and target binding, and use a target-specific restricted credential. (6) Execute once under the invocation record, and record success, denial, failure or unknown outcome.}
\req{AGT-004}{R1-B}{T-04}{No mutation between approval and dispatch may change tool version, target, resource, amount, destination or another security-relevant argument. Canonicalization is shared by authorization and execution.}
\req{AGT-005}{R1-B}{T-04}{The executor \mustnot\ validate one URL and then follow an arbitrary redirect to another destination. Address policy is resolved at connection time and credentials are bound to the registered service identity.}
\req{AGT-006}{R1-B}{T-04}{The connector uses a target-specific restricted credential, never the model's proposed credentials. State and version preconditions (ETag, version ID or transactional compare-and-swap) are passed downstream where authorization depends on target state, and a connector without that capability \mustnot\ claim state-bound mutation safety.}
\req{AGT-007}{R1-B}{T-03}{HTTP MCP tokens are validated for the gateway's audience and \mustnot\ be passed through to downstream APIs. The connector obtains its own limited downstream token; token-exchange and delegation profiles are tested rather than assumed.}
\req{AGT-008}{R1-B}{T-20}{Local stdio tools run through the trusted executor and an operating-system sandbox with bounded filesystem, process and network access. Tool servers are network-restricted so users cannot bypass the gateway with the same credentials. A remote policy service cannot intercept an unrestricted local shell or direct local file reads, and the documentation says so.}

\subsection{Approvals and reservations}
A pure policy decision can establish eligibility but cannot atomically consume a budget. Approvals, one-time grants and spending limits require stateful transactions in PostgreSQL.
\req{APR-001}{R1-B}{T-04}{A prerequisite approval or reservation is satisfied before execution. An approval binds the exact operation digest, actor and subject, target, relevant resource revisions, policy epoch, expiry and permitted uses, and a change to any material fact invalidates it.}
\req{APR-002}{R1-B}{T-05}{Approver authority and the prohibition on self-approval are checked separately. Approval consumption and reservation are atomic, and a model-generated assertion \mustnot\ stand in for the approval record.}
\req{APR-003}{R1-B}{T-05}{An approval is consumed at the documented execution checkpoint, not when it is merely displayed.}
\req{APR-004}{R1-B}{T-05}{A reservation cannot exceed its parent grant or budget under concurrency. Capacity is reserved under row-level locking or serializable transactions with unique operation and idempotency keys.}
\reqn{APR-005}{R1-B}{T-04}{The summary shown to an approver is rendered by the platform from the canonical operation plan (tool, target, amounts, destinations, resource list and highest source label). Agent- or model-supplied descriptions may be displayed only as clearly labeled unverified notes.}
\req{APR-006}{R1-B}{T-05}{Cancellation, expiry, outcome-unknown and compensation are distinct states.}
\req{APR-007}{R1-B}{T-04}{Obligations from several sources merge restrictively: field and destination sets intersect, required redactions union, and limits take the minimum. Conflicting or unsupported obligations reject, and obligations are capability-checked and enforced before their effect or disclosure.}

\subsection{Invocation state machine and unknown outcomes}
The states are PREPARED \arrowto\ RESERVED \arrowto\ DISPATCHING \arrowto\ SUCCEEDED, FAILED or UNKNOWN. A network timeout after dispatch is not proof that the external operation failed.
\req{INV-001}{R1-B}{T-05}{Invocation state transitions are recorded durably and each transition has one writer under the invocation lease.}
\req{INV-002}{R1-B}{T-05}{A unique invocation row and a transactional lease with a monotonically increasing fencing token give exactly one worker dispatch ownership. Approval consumption and budget reservation commit with the ownership transition, and in the production topology that transition passes the durability barrier of \rr{AVL-004} before the external dispatch begins.}
\req{INV-003}{R1-B}{T-05}{A replacement worker reconciles a DISPATCHING or UNKNOWN invocation and \mustnot\ redispatch it blindly.}
\req{INV-004}{R1-B}{T-05}{A retry of the same operation uses the same idempotency identity and exact argument digest. Changing arguments starts a new operation that requires new authority. Read-only calls may use a separate bounded retry policy.}
\req{INV-005}{R1-B}{T-05}{Retrying a mutation requires downstream idempotency support or target-specific reconciliation. No exactly-once promise is made across arbitrary external APIs, and tool-worker fencing cannot stop an already-running remote request unless the downstream enforces an idempotency, fencing or precondition token.}
\reqn{INV-006}{R1-B}{T-05}{An UNKNOWN invocation is resolved by connector reconciliation, or by an \code{invocation.resolve} action that records evidence and one of three outcomes: success, failure verified by evidence, or abandoned (\rr{INV-009}). The tool's maximum hold age triggers escalation and alarms (at half and at the full age) and never infers failure. For financial or destructive effects the resolver \mustnot\ be the original requester.}
\req{INV-007}{R1-B}{T-05}{A reservation for an operation with an unknown external outcome \mustnot\ be released as though no effect occurred. It remains until reconciliation or resolution commits or releases it.}
\reqn{INV-009}{R1-B}{T-05}{ABANDONED closes the workflow without establishing that the external effect did not happen. It keeps the reservation charged, keeps the approval consumed, prohibits redispatch of the operation and retains the evidence. A reservation is released only by a verified failure or by a separate \code{reservation.write-off} action, authorized to a principal other than the requester, that records the accepted risk.}
\req{INV-008}{R1-B}{T-05}{After an ambiguous failover, connectors that cannot reconcile remain UNKNOWN. The same rule applies to budget accounting and approval consumption.}

\subsection{Data flow and the context high-water mark}\label{sec:flow}
\req{FLW-001}{R1-C}{T-21}{Read permission does not constrain what a model subsequently infers, says or sends. Enforcing destination restrictions requires mediation of all relevant egress, file writes and downstream tool calls; if the model has an unrestricted independent network path, the deployment \mustnot\ claim complete output-flow enforcement.}
\req{FLW-002}{R1-C}{T-21}{A tool that reads confidential storage and sends data requires both source-read and destination-send permission, executed through one prepared plan. Raw storage and KMS credentials are not exposed to the model.}
\reqn{FLW-003}{R1-C}{T-21}{For every mediated execution context the platform records the highest effective label of any unit or derived output disclosed to it (tool results, resource reads, model context). An egress-capable call whose destination clearance is below the mark is \code{DENY}, or \code{NEEDS_APPROVAL} when a policy output says so. The mark only increases within its context.}
\reqn{FLW-004}{R1-C}{T-21}{Labels form a tenant-defined total order in the signed bundle, identified by an order generation, and the order is immutable within a generation. A unit or output with no label counts as the highest level. A mark records the generation it was computed under; when the active generation changes, a mark from an older generation is treated as the highest level unless the bundle supplies a verified monotone mapping from the old order to the new one.}
\reqn{FLW-005}{R1-C}{T-21}{The mark is conservative: it cannot tell which disclosed content influenced a call, so it over-restricts. It does not cover paths outside mediated tools, or runtimes that discard their context lifetime, and product documentation states those limits.}
\reqn{FLW-007}{R1-C}{T-21}{A mark raise is an atomic, monotonic update that is durably acknowledged before any byte of the disclosure is emitted. The decision request carries the mark revision and label-order generation, and the signed decision binds them. At egress dispatch the executor rechecks the mark revision under the invocation lease and rejects the dispatch, requiring a new decision, if the mark has risen since. Once an external call or stream is admitted and dispatched, a later raise does not retract it.}
\reqn{FLW-008}{R1-C}{T-21}{The mark belongs to a trusted execution-context lifetime. It is keyed by a context ID from a registered runtime when one is supplied, and otherwise by the delegation ID. A new MCP connection or session ID under the same key inherits the mark, and only a new key starts at the lowest level. A mark is reset early only by an explicit, audited, approval-gated \code{context.reset} action. The guarantee holds for runtimes that preserve their context lifetime (\rr{FLW-005}).}
\req{FLW-006}{R1-C}{REV}{Classification propagation and output inspection beyond the mark are separate mechanisms with explicit false-negative and false-positive limits. ``Use only for purpose X'' cannot be guaranteed by an ordinary boolean policy or by encryption alone.}

\subsection{MCP transport and surface}
\req{MCP-001}{R1-C}{T-06}{The MCP endpoint uses the official MCP Go SDK for protocol transport and JSON-RPC machinery. Project code supplies the trusted registry, policy mapping and execution boundary; a hand-written transport is not the default.}
\req{MCP-002}{R1-C}{T-06}{MCP has no independent permission rules, document store, user database or evaluator. Its handlers translate requests and responses and delegate to the shared content and decision paths.}
\req{MCP-003}{R1-C}{T-06, T-09}{Discovery and invocation are checked separately from content access. \code{tools/list} or a permit for \code{documents.read} as a tool does not authorize every document. \code{resources/read} and every content-returning tool resolve the exact resource, authorize the parent and units, and only then obtain and return permitted plaintext. A connector that cannot identify the protected content before disclosure is rejected or restricted to a fully authorized whole-resource profile, and an opaque upstream tool result is not assumed filtered.}
\req{MCP-004}{R1-C}{T-03, T-04}{Document and tool content is untrusted data, never a source of roles, permissions, trusted key locations or execution authority. Retrieved instructions cannot change the tool registry, broaden a delegation, choose a new provider endpoint or grant a follow-up tool call, and the executor checks each new effect under the same contract.}
\req{MCP-005}{R2}{T-10}{Standard MCP clients receive authorized plaintext or resource links. HPKE ciphertext is returned only to clients that explicitly support the export profile.}
\req{MCP-006}{R1-C}{T-06}{Where MCP runs in a different role process, its authenticated RPC client calls the same service implementation and preserves the represented principal and delegation.}
\reqn{MCP-007}{R1-C}{T-21}{An MCP or tool session identifier is bound to the authenticated actor, subject and delegation, and its binding revision is part of the decision request. A session identifier alone is never authority and does not by itself define the lifetime of the context mark (\rr{FLW-008}).}

\subsection{HTTP, service and future adapters}
\req{ADP-001}{R1-C}{T-06}{HTTP middleware authenticates the caller, resolves the route and resource, validates the request body and checks the exact semantic action immediately before protected work. Path parameters and request bodies are not authoritative ownership records.}
